
\vspace{0.5cm}
\noindent Confirmaremos essa equação com a aplicação do Teorema Pi. Mas, antes, observe que:
\begin{align*}
{\pi^{\prime}}_{3} &= \frac{T}{mg} \\ \\
{\pi^{\prime}}_{3} &= \frac{\frac{ML}{T^{2}}}{M \cdot \frac{L}{T^{2}}} = 1
\end{align*}

\noindent Considere a aplicação do Teorema Pi à equação, identificando as dimensões:
\begin{align*}
[T] &= \frac{ML}{T^{2}} \\\\
[g] &= \frac{L}{T^{2}} \\\\
[m] &= M \\\\
[T_{0}] &= T
\end{align*}

\noindent Portanto, a equação formatada para o grupo Pi é:
\[ {\pi^{\prime}}_{3} = (T_{0})^{\alpha} \cdot g^{\beta} \cdot m^{\delta} \cdot T \]

\noindent Substituindo as dimensões na equação, temos:
\[ {\pi^{\prime}}_{3} = T^{\alpha} \cdot \left(\frac{L}{T^{2}}\right)^{\beta} \cdot M^{\delta} \cdot \frac{ML}{T^{2}} \]

\noindent Igualando os expoentes a zero (para garantir que o grupo seja adimensional), montamos o seguinte sistema:
\[
\begin{cases}
\text{Para T: } \alpha - 2\beta - 2 = 0 \\
\text{Para L: } \beta + 1 = 0 \\
\text{Para M: } \delta + 1 = 0
\end{cases}
\]

\noindent Resolvendo o sistema, encontramos:
\[ \alpha = 0, \quad \beta = -1, \quad \delta = -1 \]

\noindent Substituindo os valores encontrados de volta nas dimensões para verificação:
\[ {\pi^{\prime}}_{3} = L^{0} \cdot \left(\frac{L}{T^{2}}\right)^{-1} \cdot M^{-1} \cdot \frac{ML}{T^{2}} = \frac{ML/T^2}{M \cdot L/T^2} = 1 \]

\noindent Assim, o grupo adimensional ${\pi^{\prime}}_{3}$ é igual a 1, confirmando que a equação (2.32) está correta sob a análise dimensional do Teorema de Buckingham Pi para a formulação revisada do problema do pêndulo.

\vspace{0.5cm}
\noindent \textbf{Conclusão:} A equação (2.32) é confirmada, pois o grupo adimensional formado é igual a 1, mostrando que a relação proposta é consistente com a análise dimensional e o Teorema de Buckingham Pi.